\subsection{Reduction of alternating bilinear forms.}

THEOREM 1. --- Let A be a principal (commutative) ring, E a free A-module of finite dimension n, and \(\Phi\) an alternating bilinear form on E. Then there exists a basis \((e_i)_{1 \leqslant i \leqslant n}\) of E and an even integer \(2r \leqslant n\) , such that

\[
1 ^ {0} \quad \Phi (e _ {1}, e _ {2}) = \alpha_ {1}, \Phi (e _ {3}, e _ {4}) = \alpha_ {2}, \dots , \Phi (e _ {2 r - 1}, e _ {2 r}) = \alpha_ {r}
\]

where the \(\alpha_{i}\) are elements \(\neq 0\) of A, and where \(\alpha_{i}\) divides \(\alpha_{i+1}\) for \(i=1,\ldots,r-1\) .

2° All the other elements \(\Phi(e_{i}, e_{j})\) where \(i \leqslant j\) are zero.

The ideals \(\mathrm{A}\alpha_{i}\) ( \(i=1,\ldots,r\) ) are uniquely determined by the preceding conditions. The submodule \(\mathrm{E}^{0}\) of E orthogonal to E is generated by \(e_{2r+1},\ldots,e_{n}\) .

We shall proceed by induction on the dimension n of E. The theorem is obvious for n = 0. If \(\Phi = 0\) , the theorem is also obvious; we may therefore assume \(\Phi \neq 0\) . Let f denote the linear map \(d_{\Phi}\) from E into \(\mathbf{E}^*\) associated on the right to \(\Phi\) ( \(\S 1, \text{no. 1}\) ); then \(f(\mathbf{E})\) is a submodule not reduced to 0 of the module \(\mathbf{E}^*\) , which is a free module of dimension \(n\) . Let \(\mathrm{A}\alpha_{1}\) be the largest invariant factor of \(f(\mathbf{E})\) with respect to \(\mathbf{E}^*\) (chap. VII, \(\S 4, \text{no. 2}\) , th. 1); we know (loc. cit.) that there exists a basis \((e_{1}', a_{2}', \ldots, a_{n}')\) of \(\mathbf{E}^*\) and an element \(f(e_{2}) \in f(\mathbf{E})\) such that \(f(e_{2}) = \alpha_{1}e_{1}'\) . Let \((e_{1}, a_{2}, \ldots, a_{n})\) the basis of E (identified with the bidual \(\mathbf{E}^{**}\) ) dual of \((e_{1}', a_{2}', \ldots, a_{n}')\) ; we have

\[
\Phi (e _ {1}, e _ {2}) = - \Phi (e _ {2}, e _ {1}) = \left\langle e _ {1}, f (e _ {2}) \right\rangle = \alpha_ {1}.\tag{1}
\]

Let P be the submodule \(Ae_{1} + Ae_{2}\) of E. We shall see that E is the direct sum of \(Ae_{1}\) , \(Ae_{2}\) and of the submodule \(P^{0}\) orthogonal to P. For this it suffices to prove that, for every \(x \in E\) , there exist elements \(\xi_{1}\) , \(\xi_{2}\) of A, uniquely determined and such that \(x - \xi_{1}e_{1} - \xi_{2}e_{2} \in P^{0}\) that is, such that

\[
\Phi (e _ {1}, x - \xi_ {1} e _ {1} - \xi_ {2} e _ {2}) = 0, \quad \Phi (e _ {2}, x - \xi_ {1} e _ {1} - \xi_ {2} e _ {2}) = 0.
\]

According to (1) these conditions are written

\[
\left\langle e _ {1}, f (x) \right\rangle = \xi_ {2} \alpha_ {1}, \left\langle e _ {2}, f (x) \right\rangle = - \xi_ {1} \alpha_ {1}.
\]

But we know (loc. cit.) that the image of \(f(E)\) under any linear form on \(E^{*}\) is contained in the ideal \(A\alpha_{1}\) , in other words all the values \(\Phi(x,y)=\langle x,f(y)\rangle\) belong to \(A\alpha_{1}\) ; whence the existence and uniqueness of \(\xi_{1}\) and \(\xi_{2}\) . Thus \(P^{0}\) is a free module of rank \(n-2\) ; there therefore exists, in \(P^{0}\) , by the induction hypothesis, a basis \((e_{3},e_{4},\ldots,e_{n})\) satisfying the conditions of the statement. To show that the basis \((e_{1},\ldots,e_{n})\) of E thus obtained also satisfies these conditions, it suffices to prove that \(\alpha_{1}\) divides \(\alpha_{2}\) ; now this follows from the fact that all the values \(\Phi(x,y)\) are multiples of \(\alpha_{1}\) . It is then clear that \(e_{2r+1},\ldots,e_{n}\) generate \(E^{0}\) . Finally, if \((e_{i}^{\prime})\) is the dual basis of \((e_{i})\) , one has \(f(e_{2j-1})=-\alpha_{j}e_{2j}^{\prime}\) and \(f(e_{2j})=\alpha_{j}e_{2j}^{\prime}\) for \(j=1,\ldots,r\) and \(f(e_{k})=0\) for \(k=2r+1,\ldots,n\) ; the ideals \(A\alpha_{1},A\alpha_{1},A\alpha_{2},A\alpha_{2},\ldots,A\alpha_{r},A\alpha_{r}\) are therefore the invariant factors of \(f(E)\) with respect to \(E^{*}\) , which proves their uniqueness (chap. VII, § 4, no. 2, th. 1).

COROLLARY 1. --- Let A be a commutative field, E a vector space of finite dimension n over A, and Φ an alternating bilinear form on E. Then there exists a basis \((e_i)_{1\leqslant i\leqslant n}\) of E and an even integer \(2r\leqslant n\) such that

\[
\Phi (\sum_ {i = 1} ^ {n} \xi_ {i} e _ {i}, \sum_ {i = 1} ^ {n} \eta_ {i} e _ {i}) = \sum_ {j = 1} ^ {r} (\xi_ {2 j - 1} \eta_ {2 j} - \xi_ {2 j} \eta_ {2 j - 1}).\tag{2}
\]

In particular \(\Phi\) is of even rank \(2r\) .

A basis satisfying (2) is called a symplectic basis for \(\Phi\) .

Remark. --- Note that this corollary is also an immediate consequence of prop. 3 of § 4, no. 2 and of its cor. 1, for with the notations of that proposition, one necessarily has H = \{0\} since Φ is alternating.

COROLLARY 2. --- Let A be a commutative field, E a vector space of finite dimension n over A. For every bivector \(z \in \bigwedge^{2} \mathrm{E}\) there exists a basis \((e_i)_{1 \leqslant i \leqslant n}\) of E such that

\[
z = e _ {1} \wedge e _ {2} + e _ {3} \wedge e _ {4} + \dots + e _ {2 r - 1} \wedge e _ {2 r} \quad (2 r \leqslant n).
\]

It suffices indeed to note that z is canonically identified with an alternating bilinear form on E* (chap. III, § 8, no. 2), and to apply cor. 1 to this form.

Translating cor. 1 into matrix language, we obtain:

COROLLARY 3. --- Let A be a commutative field, R an alternating square matrix over A. The rank of R is an even number 2r, and there exists an invertible matrix P over A such that

\[
{ } ^ { t } P . R . P = \left( \begin{array} { c c c } 0 & I _ { r } & 0 \\ - I _ { r } & 0 & 0 \\ 0 & 0 & 0 \end{array} \right) .
\]

Remark. --- If A is any commutative ring and R is an alternating square matrix of odd order n over A, then det R = 0. This follows from Cor. 3 when A is a field. A direct proof in the case where A is a field of characteristic ≠ 2 is the following: since \(^tR = -R\) , we have det \(R = \det ^t R = (-1)^n\) det R, hence 2 det \(^tR = 0\) . This being so, since the determinant of an alternating matrix ( \(\alpha_{ij}\) ) is a polynomial with integer coefficients in the \(\alpha_{ij}\) such that i \textless{} j, the principle of extension of algebraic identities (chap. IV, § 2, no. 5, Scholium) then shows that our assertion is true for an arbitrary commutative ring A.

